%%
%% part1-getting-started/ch04-book-structure.tex
%%
%% Chapter 4: Book Structure — A detailed look at the
%% frontmatter, mainmatter, and backmatter split, and how
%% each part of a MatinBook project is organized.
%%

\chapter{ساختار یک کتاب}
\label{chap:book-structure}

\section{چرا این موضوع مهم است؟}
\label{sec:book-structure-why}

در فصل قبل، یک کتاب ساده ساختیم و با پنج بخش
اصلی آن آشنا شدیم. در این فصل، نگاهی عمیق‌تر به
ساختار یک کتاب \lr{MatinBook} می‌اندازیم.

درک درست از ساختار کتاب، به شما کمک می‌کند که:

\begin{itemize}
    \item بدانید هر عنصر کتاب کجا قرار می‌گیرد.
    \item بتوانید کتاب‌های بزرگ را به‌راحتی مدیریت کنید.
    \item از خطاهای رایج جلوگیری کنید.
    \item کتاب خود را حرفه‌ای‌تر سازمان‌دهی کنید.
\end{itemize}

\mnote{اگر با ساختار \lr{LaTeX} آشنا هستید، می‌دانید که \lr{\textbackslash frontmatter}، \lr{\textbackslash mainmatter} و \lr{\textbackslash backmatter} از کلاس \lr{book} استاندارد می‌آیند.}

\section{سه بخش اصلی کتاب}
\label{sec:book-structure-three-parts}

هر کتاب \lr{MatinBook} از سه بخش اصلی تشکیل شده
است که با سه دستور مشخص می‌شوند:

\begin{enumerate}
    \item \textbf{پیش‌متن} (\lr{\textbackslash frontmatter}):
    بخش‌هایی که پیش از متن اصلی می‌آیند.
    \item \textbf{متن اصلی} (\lr{\textbackslash mainmatter}):
    فصل‌های اصلی کتاب.
    \item \textbf{پس‌متن} (\lr{\textbackslash backmatter}):
    بخش‌هایی که پس از متن اصلی می‌آیند.
\end{enumerate}

\mnote{هر یک از این سه دستور، شماره‌گذاری صفحات را تغییر می‌دهد: پیش‌متن با حروف ابجد، متن اصلی با اعداد عربی، پس‌متن بدون شماره.}

\section{پیش‌متن}
\label{sec:book-structure-frontmatter}

پیش‌متن، بخش‌هایی است که پیش از فصل‌های اصلی
کتاب قرار می‌گیرند. این بخش‌ها شامل موارد زیر
هستند:

\begin{itemize}
    \item \textbf{صفحه‌ی عنوان:} با \cmd{maketitle}.
    \item \textbf{شناسنامه:} صفحه‌ی حقوق و
    اطلاعات نشر.
    \item \textbf{تقدیم‌نامه:} صفحه‌ی تقدیم.
    \item \textbf{پیشگفتار:} متن مقدمه.
    \item \textbf{فهرست مطالب:} با
    \cmd{tableofcontents}.
    \item \textbf{فهرست تصاویر:} با
    \cmd{listoffigures}.
    \item \textbf{فهرست جداول:} با
    \cmd{listoftables}.
    \item \textbf{فهرست الگوریتم‌ها:} با
    \cmd{listofalgorithms}.
\end{itemize}

\subsection{شماره‌گذاری صفحات در پیش‌متن}
\label{sec:book-structure-frontmatter-numbering}

در پیش‌متن، شماره‌گذاری صفحات با \textbf{حروف
ابجد فارسی} انجام می‌شود: الف، ب، پ، ت، ث، ج،
چ، ح، خ، د، ذ، ر، ز، ژ، س، ش، ص، ض، ط، ظ، ع،
غ، ف، ق، ک، گ، ل، م، ن، و، ه، ی.

این کار، مطابق استاندارد نشر ایران است و باعث
می‌شود که پیش‌متن، مستقل از متن اصلی شماره‌گذاری
شود.

\mnote{اگر پیش‌متن طولانی باشد، ممکن است حروف ابجد تمام شوند. \lr{xepersian} به‌طور خودکار از ترکیب‌های جدید استفاده می‌کند.}

\subsection{هندسه‌ی صفحه در پیش‌متن}
\label{sec:book-structure-frontmatter-geometry}

در پیش‌متن، حاشیه‌ی بیرونی پهن وجود ندارد و
صفحه، حالت \textbf{متقارن} دارد:

\begin{itemize}
    \item حاشیه‌ی داخلی: \pnum{۱.۵} سانتی‌متر.
    \item حاشیه‌ی بیرونی: \pnum{۱.۵} سانتی‌متر.
    \item عرض متن: حدود \pnum{۱۴.۳} سانتی‌متر.
\end{itemize}

برای فعال‌سازی این حالت، از دستور
\cmd{frontmattergeometry} استفاده کنید.

\mnote{دلیل حذف حاشیه‌ی پهن در پیش‌متن این است که این بخش‌ها به یادداشت حاشیه‌ای نیاز ندارند و می‌توانند از عرض کامل متن استفاده کنند.}

\subsection{ساختار یک پیش‌متن نمونه}
\label{sec:book-structure-frontmatter-example}

یک پیش‌متن نمونه در \lr{MatinBook} به این صورت است:

\begin{latin}
\begin{minted}{latex}
\frontmatter
\frontmattergeometry

% |\pc{صفحه‌ی عنوان}|
\maketitle

% |\pc{شناسنامه}|
\thispagestyle{empty}
\vfill
\begin{center}
    |\pc{تمامی حقوق محفوظ است.}|
\end{center}
\clearpage

% |\pc{تقدیم‌نامه}|
\thispagestyle{empty}
\vfill
\begin{flushright}
    |\pc{تقدیم به ...}|
\end{flushright}
\clearpage

% |\pc{پیشگفتار}|
\chapter*{|\pc{پیشگفتار}|}
\addcontentsline{toc}{chapter}{|\pc{پیشگفتار}|}
|\pc{متن پیشگفتار...}|

% |\pc{فهرست‌ها}|
\tableofcontents
\clearpage
\listoffigures
\clearpage
\listoftables
\clearpage
\end{minted}
\end{latin}

\mnote{دستور \lr{\textbackslash chapter*} یک فصل بدون شماره می‌سازد. برای پیشگفتار و پیوست‌های بدون شماره مناسب است.}

\subsection{دستور \cmd{chapter*} در پیش‌متن}
\label{sec:book-structure-chapter-star}

برای بخش‌هایی مانند پیشگفتار که نباید شماره‌ی
فصل داشته باشند، از دستور \cmd{chapter*} استفاده
می‌کنیم. این دستور، یک عنوان بدون شماره می‌سازد.

برای اینکه این بخش در فهرست مطالب هم بیاید،
باید دستور \cmd{addcontentsline} را نیز اضافه
کنیم:

\begin{latin}
\begin{minted}{latex}
\chapter*{|\pc{پیشگفتار}|}
\addcontentsline{toc}{chapter}{|\pc{پیشگفتار}|}
\end{minted}
\end{latin}

\mnote{دستور \lr{\textbackslash addcontentsline} سه آرگومان می‌گیرد: فایل فهرست (\lr{toc})، سطح (\lr{chapter})، و متن.}

\section{متن اصلی}
\label{sec:book-structure-mainmatter}

متن اصلی، بخش اصلی کتاب است و شامل فصل‌های
اصلی می‌شود. این بخش با دستور \cmd{mainmatter}
شروع می‌شود.

\subsection{شماره‌گذاری صفحات در متن اصلی}
\label{sec:book-structure-mainmatter-numbering}

در متن اصلی، شماره‌گذاری صفحات با \textbf{اعداد
عربی} انجام می‌شود: ۱، ۲، ۳، ...

\subsection{هندسه‌ی صفحه در متن اصلی}
\label{sec:book-structure-mainmatter-geometry}

در متن اصلی، حاشیه‌ی بیرونی پهن (به‌طور پیش‌فرض
\pnum{۵.۹} سانتی‌متر) فعال می‌شود:

\begin{itemize}
    \item حاشیه‌ی داخلی: \pnum{۱.۵} سانتی‌متر.
    \item حاشیه‌ی بیرونی: \pnum{۵.۹} سانتی‌متر.
    \item عرض متن: حدود \pnum{۱۰.۳} سانتی‌متر.
\end{itemize}

برای فعال‌سازی این حالت، از دستور
\cmd{mainmattergeometry} استفاده کنید.

\mnote{حاشیه‌ی پهن برای یادداشت‌های حاشیه‌ای، شکل‌های کوچک و فرمول‌های مکمل طراحی شده است.}

\subsection{ساختار یک فصل نمونه}
\label{sec:book-structure-chapter-example}

یک فصل نمونه در \lr{MatinBook} به این صورت است:

\begin{latin}
\begin{minted}{latex}
\chapter{|\pc{فصل اول}|}
\label{chap:first}

\section{|\pc{مقدمه}|}
\label{sec:first-intro}

|\pc{متن مقدمه...}|

\section{|\pc{مفهوم اصلی}|}
\label{sec:first-main}

|\pc{متن بخش اصلی...}|

\begin{theorem}[|\pc{قضیه‌ی نمونه}|]
    |\pc{متن قضیه...}|
    \label{thm:first}
\end{theorem}

\begin{proof}
    |\pc{اثبات قضیه...}|
\end{proof}

\section{|\pc{نتیجه‌گیری}|}
\label{sec:first-conclusion}

|\pc{متن نتیجه‌گیری...}|
\end{minted}
\end{latin}

\mnote{هر فصل باید حداقل یک \lr{\textbackslash label} داشته باشد تا بتوان به آن ارجاع داد.}

\section{پس‌متن}
\label{sec:book-structure-backmatter}

پس‌متن، بخش‌هایی است که پس از فصل‌های اصلی
کتاب قرار می‌گیرند. این بخش‌ها شامل موارد زیر
هستند:

\begin{itemize}
    \item \textbf{منابع و مراجع:} با
    \cmd{printbibliography}.
    \item \textbf{نمایه:} با \cmd{printindex}.
    \item \textbf{پیوست‌ها:} با \cmd{appendix}.
\end{itemize}

\subsection{شماره‌گذاری صفحات در پس‌متن}
\label{sec:book-structure-backmatter-numbering}

در پس‌متن، شماره‌گذاری صفحات همانند متن اصلی
با اعداد عربی ادامه می‌یابد.

\subsection{هندسه‌ی صفحه در پس‌متن}
\label{sec:book-structure-backmatter-geometry}

در پس‌متن، هندسه‌ی صفحه به حالت \textbf{متقارن}
برمی‌گردد (همانند پیش‌متن). برای این کار، از
دستور \cmd{frontmattergeometry} استفاده می‌کنیم.

\mnote{استفاده از \lr{\textbackslash frontmattergeometry} در پس‌متن، کمی گیج‌کننده است. این دستور فقط «هندسه‌ی متقارن» را فعال می‌کند، نه «حالت پیش‌متن» را.}

\subsection{پیوست‌ها}
\label{sec:book-structure-appendices}

پیوست‌ها بخش‌هایی هستند که در انتهای کتاب
می‌آیند و معمولاً شامل مطالب تکمیلی، جدول‌ها،
یا کدهای طولانی هستند. برای شروع پیوست‌ها،
از دستور \cmd{appendix} استفاده می‌کنیم:

\begin{latin}
\begin{minted}{latex}
\appendix

\chapter{|\pc{پیوست الف}|}
\label{app:first}

|\pc{متن پیوست...}|
\end{minted}
\end{latin}

پس از \cmd{appendix}، شماره‌گذاری فصل‌ها به
حروف ابجد لاتین (A, B, C, ...) تغییر می‌کند.

\mnote{پیوست‌ها معمولاً در \lr{\textbackslash mainmatter} قرار می‌گیرند، نه در \lr{\textbackslash backmatter}.}

\subsection{ساختار یک پس‌متن نمونه}
\label{sec:book-structure-backmatter-example}

یک پس‌متن نمونه در \lr{MatinBook} به این صورت
است:

\begin{latin}
\begin{minted}{latex}
% |\pc{پیوست‌ها}| (|\pc{در متن اصلی}|)
\appendix

\chapter{|\pc{پیوست الف}|}
\label{app:a}

|\pc{متن پیوست...}|

% |\pc{پس‌متن}| (|\pc{بدون شماره}|)
\backmatter
\frontmattergeometry

% |\pc{منابع}| (|\pc{بدون}|\lr{latin})
\printbibliography[title={|\pc{منابع و مراجع}|}]

% |\pc{نمایه}|
\printindex
\end{minted}
\end{latin}

\mnote{دستور \lr{\textbackslash printbibliography} را داخل \lr{latin} قرار ندهید. \lr{xepersian} جهت متن را به‌درستی مدیریت می‌کند.}

\section{جدول مقایسه سه بخش}
\label{sec:book-structure-comparison}

جدول \cref{tab:book-structure} مقایسه‌ی خلاصه‌ای
بین سه بخش اصلی کتاب را نشان می‌دهد.

\begin{matintable}{مقایسه‌ی سه بخش اصلی کتاب}{tab:book-structure}
\begin{tblr}{
    width = \textwidth,
    colspec = {X[l] X[c] X[c] X[c]},
    hlines, vlines,
    colsep = 6pt,
    rowsep = 4pt,
    row{1} = {font=\bfseries},
}
ویژگی & پیش‌متن & متن اصلی & پس‌متن \\
دستور شروع & \lr{\textbackslash frontmatter} & \lr{\textbackslash mainmatter} & \lr{\textbackslash backmatter} \\
شماره‌گذاری & حروف ابجد & اعداد عربی & اعداد عربی \\
هندسه & متقارن & حاشیه‌ی پهن & متقارن \\
محتوا & پیشگفتار، فهرست & فصل‌ها & منابع، نمایه \\
شماره‌ی فصل & ندارد & دارد & ندارد \\
\end{tblr}
\end{matintable}

\section{ترتیب دقیق عناصر}
\label{sec:book-structure-order}

ترتیب دقیق عناصر کتاب در \lr{MatinBook} به
این صورت است:

\begin{enumerate}
    \item \textbf{جلد جلو:} \cmd{makecover}.
    \item \textbf{پیش‌متن:} \cmd{frontmatter} +
    \cmd{frontmattergeometry}.
    \begin{enumerate}
        \item صفحه‌ی عنوان: \cmd{maketitle}.
        \item شناسنامه.
        \item تقدیم‌نامه.
        \item پیشگفتار: \cmd{chapter*} +
        \cmd{addcontentsline}.
        \item فهرست مطالب: \cmd{tableofcontents}.
        \item فهرست تصاویر: \cmd{listoffigures}.
        \item فهرست جداول: \cmd{listoftables}.
        \item فهرست الگوریتم‌ها: \cmd{listofalgorithms}.
    \end{enumerate}
    \item \textbf{متن اصلی:} \cmd{mainmatter} +
    \cmd{mainmattergeometry}.
    \begin{enumerate}
        \item بخش‌ها: \cmd{part}.
        \item فصل‌ها: \cmd{chapter}.
        \item پیوست‌ها: \cmd{appendix} +
        \cmd{chapter}.
    \end{enumerate}
    \item \textbf{پس‌متن:} \cmd{backmatter} +
    \cmd{frontmattergeometry}.
    \begin{enumerate}
        \item منابع: \cmd{printbibliography}.
        \item نمایه: \cmd{printindex}.
    \end{enumerate}
    \item \textbf{جلد عقب:} \cmd{makebackcover}.
\end{enumerate}

\mnote{ترتیب دقیق این عناصر، یکی از مهم‌ترین عوامل در تولید کتاب حرفه‌ای است.}

\section{خطاهای رایج}
\label{sec:book-structure-mistakes}

\subsection{فراموش کردن \cmd{frontmattergeometry}}
\label{sec:book-structure-mistake-frontgeo}

اگر \cmd{frontmattergeometry} را فراموش کنید،
پیش‌متن با حاشیه‌ی پهن نمایش داده می‌شود که
برای فهرست مطالب و پیشگفتار مناسب نیست.

\textbf{راه‌حل:} همیشه \cmd{frontmattergeometry} را
پس از \cmd{frontmatter} قرار دهید.

\subsection{فراموش کردن \cmd{mainmattergeometry}}
\label{sec:book-structure-mistake-maingeo}

اگر \cmd{mainmattergeometry} را فراموش کنید،
متن اصلی با حاشیه‌ی متقارن (بدون حاشیه‌ی پهن)
نمایش داده می‌شود که برای یادداشت‌های حاشیه‌ای
مشکل ایجاد می‌کند.

\textbf{راه‌حل:} همیشه \cmd{mainmattergeometry} را
پس از \cmd{mainmatter} قرار دهید.

\subsection{قرار دادن پیوست‌ها در پس‌متن}
\label{sec:book-structure-mistake-appendix}

پیوست‌ها باید در متن اصلی قرار گیرند، نه در
پس‌متن. اگر پیوست‌ها را در پس‌متن قرار دهید،
شماره‌گذاری آن‌ها به‌درستی انجام نمی‌شود.

\textbf{راه‌حل:} پیوست‌ها را پس از فصل‌های اصلی
و پیش از \cmd{backmatter} قرار دهید.

\subsection{فراموش کردن \cmd{addcontentsline} برای \cmd{chapter*}}
\label{sec:book-structure-mistake-addcontentsline}

اگر از \cmd{chapter*} برای پیشگفتار استفاده
کنید و \cmd{addcontentsline} را فراموش کنید،
پیشگفتار در فهرست مطالب ظاهر نمی‌شود.

\textbf{راه‌حل:} همیشه \cmd{addcontentsline} را
پس از \cmd{chapter*} قرار دهید.

\section{تمرین‌ها}
\label{sec:book-structure-exercises}

\begin{exercise}
    یک کتاب با پیش‌متن کامل (شناسنامه، تقدیم،
    پیشگفتار، فهرست مطالب) بسازید. آیا
    شماره‌گذاری صفحات با حروف ابجد انجام
    می‌شود؟
    \label{exc:book-structure-frontmatter}
\end{exercise}

\begin{exercise}
    یک کتاب با دو بخش و چهار فصل بسازید.
    از دستور \cmd{part} برای ایجاد بخش‌ها
    استفاده کنید.
    \label{exc:book-structure-parts}
\end{exercise}

\begin{exercise}
    یک پیوست به کتاب خود اضافه کنید و
    شماره‌گذاری آن را بررسی کنید.
    \label{exc:book-structure-appendix}
\end{exercise}

\begin{exercise}
    جدول \cref{tab:book-structure} را با
    کتاب خود مقایسه کنید. آیا همه‌ی
    ویژگی‌ها را رعایت کرده‌اید؟
    \label{exc:book-structure-compare}
\end{exercise}

\section{خلاصه}
\label{sec:book-structure-summary}

در این فصل، ساختار یک کتاب \lr{MatinBook} را
به‌طور مفصل بررسی کردیم:

\begin{itemize}
    \item یک کتاب از سه بخش اصلی تشکیل شده:
    پیش‌متن، متن اصلی، پس‌متن.

    \item پیش‌متن با حروف ابجد فارسی و هندسه‌ی
    متقارن شماره‌گذاری می‌شود.

    \item متن اصلی با اعداد عربی و حاشیه‌ی پهن
    شماره‌گذاری می‌شود.

    \item پس‌متن با اعداد عربی و هندسه‌ی متقارن
    شماره‌گذاری می‌شود.

    \item پیوست‌ها در متن اصلی قرار می‌گیرند،
    نه در پس‌متن.

    \item برای بخش‌های بدون شماره، از
    \cmd{chapter*} همراه با \cmd{addcontentsline}
    استفاده می‌شود.

    \item خطاهای رایج شامل فراموش کردن
    \cmd{frontmattergeometry}، \cmd{mainmattergeometry}،
    قرار دادن پیوست‌ها در پس‌متن، و فراموش
    کردن \cmd{addcontentsline} است.
\end{itemize}

این فصل، آخرین فصل از بخش اول کتاب (شروع به
کار) بود. در بخش دوم (اجزای کتاب)، هر یک از
اجزای کتاب را به‌طور مفصل بررسی خواهیم کرد.